\exercisesection{\S~1}

We denote by \(\mathfrak{s}\) the Lie algebra \(\mathfrak{sl}(2,k)\) .

\begin{enumerate}
\def\labelenumi{\arabic{enumi})}
\tightlist
\item
  Let \(\mathrm{U}\) be the enveloping algebra of \(\mathfrak{s}\) . Show that the element
\end{enumerate}

\[
\mathrm{C} = H ^ {2} - 2 (X _ {+} X _ {-} + X _ {-} X _ {+}) = H ^ {2} + 2 H - 4 X _ {-} X _ {+}
\]

belongs to the centre of U, and that its image in the representation associated to \(V(m)\) is the homothety with ratio \(m(m+2)\) .

\begin{enumerate}
\def\labelenumi{\arabic{enumi})}
\setcounter{enumi}{1}
\tightlist
\item
  Let \(\lambda \in k\) , let \(Z(\lambda)\) be a vector space having a basis \((e_n)\) , with \(n = 0,1,\ldots\) , and let \(X_{+}, X_{-}, H\) be the endomorphisms of \(Z(\lambda)\) defined by formulas (2) of Prop. 1.
\end{enumerate}

\begin{enumerate}
\def\labelenumi{\alph{enumi})}
\item
  Verify that this gives an \(\mathfrak{s}\) -module structure on \(Z(\lambda)\) .
\item
  Assume that \(\lambda\) is not an integer \(\geq 0\) . Show that the \(\mathfrak{s}\) -module \(Z(\lambda)\) is simple.
\item
  Assume that \(\lambda\) is an integer \(\geq 0\) . Let \(Z'\) be the subspace of \(Z(\lambda)\) generated by the \(e_{n}\) , \(n > \lambda\) . Show that \(Z'\) is an s-submodule of \(Z(\lambda)\) , isomorphic to \(Z(-\lambda - 2)\) , and that the quotient \(Z(\lambda)/Z'\) is isomorphic to the simple s-module \(V(\lambda)\) . Show that the only s-submodules of \(Z(\lambda)\) are 0, \(Z'\) and \(Z(\lambda)\) .
\end{enumerate}

\begin{enumerate}
\def\labelenumi{\arabic{enumi})}
\setcounter{enumi}{2}
\tightlist
\item
  Let \(\mathrm{E}\) be a vector space having a basis \((e_n)_{n\in \mathbf{Z}}\) . Let
\end{enumerate}

\[
a (n) = a _ {0} + a _ {1} n, b (n) = b _ {0} + b _ {1} n, c (n) = c _ {0} + c _ {1} n
\]

be three affine functions with coefficients in k. Define endomorphisms \(X_{+}, X_{-}\) , H of E by the formulas

\[
X _ {+} e _ {n} = a (n) e _ {n - 1}, X _ {-} e _ {n} = b (n) e _ {n + 1}, H e _ {n} = c (n) e _ {n}.
\]

This gives an \(\mathfrak{s}\) -module structure on E if and only if:

\[
a _ {1} b _ {1} = 1, c _ {1} = - 2, c _ {0} = - a _ {0} b _ {1} - a _ {1} b _ {0}.
\]

Under what condition is this \(\mathfrak{s}\) -module simple?

¶ 4) Let g be a Lie algebra. A g-module E is said to be locally finite if it is a union of finite dimensional g-submodules; equivalently, every g-submodule of E of finite type (as a U(g)-module) is finite dimensional.

\begin{enumerate}
\def\labelenumi{\alph{enumi})}
\item
  Let \(0 \rightarrow E' \rightarrow E \rightarrow E'' \rightarrow 0\) be an exact sequence of g-modules. Show that, if \(E'\) and \(E''\) are locally finite, so is E (reduce to the case in which E is of finite type, and use the fact that \(\mathrm{U}(\mathfrak{g})\) is noetherian, cf.~Chap. I, §2, no. 6).
\item
  Assume that g is semi-simple. Show that E is locally finite if and only if E is a direct sum of finite dimensional simple g-modules.
\item
  Assume that \(\mathfrak{g} = \mathfrak{s}\) . Show that E is locally finite if and only if the following conditions are satisfied:
\end{enumerate}

\(c_{1})\) E has a basis consisting of eigenvectors of H,

\(c_{2}\) ) the endomorphisms \(X_{+\mathrm{E}}\) and \(X_{-\mathrm{E}}\) are locally nilpotent.

(Begin by showing that, if E satisfies \(c_{1}\) ) and \(c_{2}\) ), and is not reduced to 0, it contains a primitive element e of integer weight \(m \geq 0\) ; prove next that \(X^{m+1}e = 0\) , and hence that E contains \(V(m)\) . Conclude by applying a) to the largest locally finite submodule \(E'\) of E.)

\begin{enumerate}
\def\labelenumi{\arabic{enumi})}
\setcounter{enumi}{4}
\tightlist
\item
  Let E be a locally finite \(\mathfrak{s}\) -module (Exerc. 4). For all \(m \geq 0\) , denote by \(\mathrm{L}(m)\) the vector space of \(\mathfrak{s}\) -homomorphisms from \(\mathrm{V}(m)\) to E.
\end{enumerate}

\begin{enumerate}
\def\labelenumi{\alph{enumi})}
\item
  Define an isomorphism from E to \(\bigoplus_{m} \mathrm{L}_{m} \otimes \mathrm{V}(m)\) .
\item
  Denote by \(\Phi_{m}\) the invariant bilinear form on \(\mathrm{V}(m)\) defined in no. 3, Remark 3. For all \(m \in N\) , let \(b_{m}\) be a bilinear form on \(L_{m}\) ; let b be the bilinear form on E which corresponds, via the isomorphism in a), to the direct sum of the forms \(b_{m} \otimes \Phi_{m}\) . Show that b is invariant, and that every invariant bilinear form on E is obtained in this way in a unique manner. The form b is symmetric (resp. alternating) if and only if the \(b_{m}\) , m even, are symmetric (resp. alternating), and the \(b_{m}\) , m odd, are alternating (resp. symmetric). The form b is non-degenerate if and only if the \(b_{m}\) are non-degenerate.
\item
  Assume that E is finite dimensional. Show that E is monogenic (as a \(\mathrm{U}(\mathfrak{s})\) -module) if and only if \(\dim L_{m} \leq m + 1\) for all \(m \geq 0\) .
\end{enumerate}

\begin{enumerate}
\def\labelenumi{\arabic{enumi})}
\setcounter{enumi}{5}
\tightlist
\item
  If E is a finite dimensional \(\mathfrak{s}\) -module, and \(n\) an integer, denote by \(a_{n}\) the dimension of the eigenspace of \(H_{\mathrm{E}}\) relative to the eigenvalue \(n\) . Denote by \(c_{\mathrm{E}}(\mathrm{T})\) the element of \(\mathbf{Z}[\mathrm{T},\mathrm{T}^{-1}]\) defined by \(c_{\mathrm{E}}(\mathrm{T}) = \sum_{n\in \mathbf{Z}}a_n\mathrm{T}^n\) .
\end{enumerate}

\begin{enumerate}
\def\labelenumi{\alph{enumi})}
\tightlist
\item
  Define \(\mathrm{L}_m\) as in Exerc. 5. Show that
\end{enumerate}

\(\dim \mathrm{L}_m = a_m - a_{m + 2}\) for every integer \(m\geq 0\)

Deduce that \(c_{\mathrm{E}}(\mathrm{T}) = c_{\mathrm{E}'}(\mathrm{T})\) if and only if E and \(E'\) are isomorphic. Recover this result by using Exerc. 18 e) of Chap. VII, §3.

\begin{enumerate}
\def\labelenumi{\alph{enumi})}
\setcounter{enumi}{1}
\item
  Show that \(c_{E\oplus F} = c_{E} + c_{F}\) and \(c_{E\otimes F} = c_{E}.c_{F}\) .
\item
  We have \(c_{\mathrm{V}(m)}(\mathrm{T}) = (\mathrm{T}^{m + 1} - \mathrm{T}^{-m - 1}) / (\mathrm{T} - \mathrm{T}^{-1})\) .
\end{enumerate}

\(d)\) Deduce from \(a\) , \(b\) , \(c)\) that, if \(m \geq m' \geq 0\) , the \(\mathfrak{s}\) -module \(\mathrm{V}(m) \otimes \mathrm{V}(m')\) is isomorphic to

\[
\mathrm{V} (m + m ^ {\prime}) \oplus \mathrm{V} (m + m ^ {\prime} - 2) \oplus \mathrm{V} (m + m ^ {\prime} - 4) \oplus \dots \oplus \mathrm{V} (m - m ^ {\prime}).
\]

If \(m\) is \(\geq 0\) , the \(\mathfrak{s}\) -module \(\mathbf{S}^2\mathrm{V}(m)\) is isomorphic to

\[
\mathrm{V} (2 m) \oplus \mathrm{V} (2 m - 4) \oplus \mathrm{V} (2 m - 8) \oplus \dots \oplus \left\{ \begin{array}{l l} \mathrm{V} (0) & \text {if m is even} \\ \mathrm{V} (2) & \text {if m is odd} \end{array} \right.
\]

and the \(\mathfrak{s}\) -module \(\bigwedge^2\mathrm{V}(m)\) is isomorphic to

\[
\mathrm{V} (2 m - 2) \oplus \mathrm{V} (2 m - 6) \oplus \mathrm{V} (2 m - 10) \oplus \dots \oplus \left\{ \begin{array}{l l} \mathrm{V} (2) & \text {if m is even} \\ \mathrm{V} (0) & \text {if m is odd.} \end{array} \right.
\]

\begin{enumerate}
\def\labelenumi{\arabic{enumi})}
\setcounter{enumi}{6}
\item
  Let E be a finite dimensional s-module. Show that the dual of E is isomorphic to E.
\item
  Show that the \(\mathfrak{s}\) -module \(V(m)\) can be realized as the space of homogeneous polynomials \(f(u,v)\) of degree \(m\) in two variables, the operators \(X_{+}, X_{-}\) and \(H\) being given by:
\end{enumerate}

\[
X _ {+} f = u \frac {\partial f}{\partial v}, X _ {-} f = - v \frac {\partial f}{\partial u}, H f = u \frac {\partial f}{\partial u} - v \frac {\partial f}{\partial v}.
\]

¶9) Consider the operation of \(\mathfrak{s}\) , via the adjoint representation, on its enveloping algebra U and on its symmetric algebra \(\mathbf{S} = \bigoplus \mathbf{S}^n\) .

\begin{enumerate}
\def\labelenumi{\alph{enumi})}
\tightlist
\item
  Determine the weights of the \(\mathfrak{s}\) -module \(\mathbf{S}^n\) . Deduce the following isomorphisms of \(\mathfrak{s}\) -modules:
\end{enumerate}

\[
\begin{array}{l l} \mathbf {S} ^ {n} \longrightarrow \mathrm{V} (2 n) \oplus \mathrm{V} (2 n - 4) \oplus \mathrm{V} (2 n - 8) \oplus \dots \oplus \mathrm{V} (0) & n \text {even} \\ \mathbf {S} ^ {n} \longrightarrow \mathrm{V} (2 n) \oplus \mathrm{V} (2 n - 4) \oplus \mathrm{V} (2 n - 8) \oplus \dots \oplus \mathrm{V} (2) & n \text {odd}. \end{array}
\]

In particular, the elements of \(\mathbf{S}^n\) invariant under \(\mathfrak{s}\) form a space of dimension 1 (resp. 0) if \(n\) is even (resp. odd).

\begin{enumerate}
\def\labelenumi{\alph{enumi})}
\setcounter{enumi}{1}
\item
  Show that the subalgebra of S consisting of the elements invariant under s is the polynomial algebra \(k[\Gamma]\) , where \(\Gamma = H^{2} - 4X_{-}X_{+}\) . The module \(S^{n}/\Gamma.S^{n-2}\) is isomorphic to \(\mathrm{V}(2n)\) .
\item
  Show that the centre of U is the polynomial algebra \(k[\mathbf{C}]\) , where C is the element defined in Exerc. 1 (use b) and the Poincaré-Birkhoff-Witt theorem).
\end{enumerate}

\begin{enumerate}
\def\labelenumi{\arabic{enumi})}
\setcounter{enumi}{9}
\tightlist
\item
  Let \(m\) be an integer \(> 0\) , \(\mathbf{S}\) the graded algebra \(k[\mathrm{X}_1, \ldots, \mathrm{X}_m]\) , and \(a_1, \ldots, a_m\) elements of \(k^*\) . Let \(\Phi = \sum_{i,j=1}^{m} a_{ij} X_i X_j\) be a quadratic form; put \(\mathrm{D}_i = \frac{\partial}{\partial X_i}\) and \(\mathrm{D} = \sum_{i,j=1}^{m} b_{ij} \mathrm{D}_i \mathrm{D}_j\) where \((b_{ij})\) is the inverse matrix of \((a_{ij})\) .
\end{enumerate}

\begin{enumerate}
\def\labelenumi{\alph{enumi})}
\tightlist
\item
  Show that there exists a unique \(\mathfrak{s}\) -module structure on \(\mathbf{S}\) such that \(X_{+}f = \frac{1}{2}\mathrm{D}(f), X_{-}f = \frac{1}{2}\Phi f\) and \(Hf = -(m / 2 + n)f\) if \(f\) is homogeneous of degree \(n\) . b) Put \(\mathrm{A}^n = \mathbf{S}^n \cap \mathrm{Ker}(\mathrm{D})\) . Show that \(\mathbf{S}^n\) is the direct sum of the subspaces \(\Phi^p\) . \(\mathrm{A}^q\) for \(2p + q = n\) . Deduce the identity
\end{enumerate}

\[
\sum_ {n = 0} ^ {\infty} \dim (\mathrm{A} ^ {n}) \mathrm{T} ^ {n} = (1 + \mathrm{T}) / (1 - \mathrm{T}) ^ {m - 1}.
\]

\begin{enumerate}
\def\labelenumi{\alph{enumi})}
\setcounter{enumi}{2}
\item
  If \(f\) is a non-zero element of \(\mathbf{A}^n\) , the \(\varPhi^p f\) , \(p \geq 0\) , form a basis of a simple \(\mathfrak{s}\) -submodule of \(\mathbf{S}\) , isomorphic to the module \(Z\left(-\frac{m}{2} - n\right)\) of Exerc. 2.
\item
  Make the cases \(m = 1\) and \(m = 2\) explicit. Use the case \(m = 3\) to recover the results of Exerc. 9.
\end{enumerate}

\begin{enumerate}
\def\labelenumi{\arabic{enumi})}
\setcounter{enumi}{10}
\tightlist
\item
  Assume that \(k\) is algebraically closed. Let \(\mathfrak{g}\) be a Lie algebra, \(V\) a simple \(\mathfrak{g}\) -module, and \(D\) the field of endomorphisms of the \(\mathfrak{g}\) -module \(V\) . Show that, if \(k\) is uncountable, \(^7\) then \(D = k\) . (If not, \(D\) would contain a subfield isomorphic to \(k(X)\) , \(X\) being an indeterminate, and we would have \(\dim_k D > \aleph_0\) , hence also \(\dim_k V > \aleph_0\) , which is absurd since \(V\) is a monogenic \(U(\mathfrak{g})\) -module.)
\end{enumerate}

\footnotetext[7]{The conclusion remains true even if k is countable, cf. D. QUILLEN, On the endomorphism ring of a simple module over an enveloping algebra, Proc. Amer. Math. Soc., Vol. XXI (1969), pp. 171-172.}

¶ 12) Let \(q \in k\) , and let \(W\) be a vector space with basis \((e_0, e_1, e_2, \ldots)\) .

\begin{enumerate}
\def\labelenumi{\alph{enumi})}
\tightlist
\item
  Show that there exists a unique representation \(\rho_{q}\) of \(\mathfrak{s}\) on W such that
\end{enumerate}

\[
\begin{array}{r l} & {\rho_ {q} (H) e _ {n} = 2 e _ {n + 1}, \quad \rho_ {q} (X _ {+}) e _ {n} = (\frac {1}{2} \rho_ {q} (H) - 1) ^ {n} e _ {0},} \\ & {\rho_ {q} (X _ {-}) e _ {n} = (\frac {1}{2} \rho_ {q} (H) + 1) ^ {n} (- q e _ {0} + e _ {1} + e _ {2}).} \end{array}
\]

We have \(\rho_{q}(\mathrm{C}) = 4q\) , where \(C = H^{2} + 2H - 4X_{-}X_{+}\) (cf.~Exerc. 1). The representation \(\rho_{q}\) is simple. The elements \(x \in s\) such that \(\rho_{q}(x)\) admits an eigenvalue are the multiples of \(X_{+}\) . The endomorphism \(\rho_{q}(X_{+})\) admits 1 as an eigenvalue, with multiplicity 1.

\begin{enumerate}
\def\labelenumi{\alph{enumi})}
\setcounter{enumi}{1}
\tightlist
\item
  Let \(\rho\) be a simple representation of \(\mathfrak{s}\) such that \(\rho(\mathrm{C}) = 4q\) and such that \(\rho(X_{+})\) has 1 as an eigenvalue. Show that \(\rho\) is equivalent to \(\rho_q\) .
\end{enumerate}

¶ 13) Assume that \(k = \mathbf{C}\) . Put \(C = H^2 + 2H - 4X_-X_+\) , cf.~Exerc. 1. A representation \(\rho\) of \(\mathfrak{s} = \mathfrak{sl}(2, \mathbf{C})\) is said to be \(H\) -diagonalizable if the underlying space of \(\rho\) has a basis consisting of eigenvectors of \(\rho(H)\) . Let \(q \in \mathbf{C}\) and \(v \in \mathbf{Q}/\mathbf{Z}\) .

\begin{enumerate}
\def\labelenumi{\alph{enumi})}
\tightlist
\item
  Let S be a vector space with a basis \((e_{w})_{w\in\mathbf{C}}\) indexed by the elements of C. Let \(S_{v}=\sum_{w\in v}Ce_{w}\) . There exists a unique representation \(\rho_{v,q}\) of s on \(S_{v}\) such that
\end{enumerate}

\[
\rho_ {v, q} (X _ {+}) e _ {w} = (q - w ^ {2} - w) ^ {1 / 2} e _ {w + 1}
\]

\[
\begin{array}{c} \rho_ {v, q} (X _ {-}) e _ {w} = (q - w ^ {2} + w) ^ {1 / 2} e _ {w - 1} \\ \rho_ {v, q} (H) e _ {w} = 2 w e _ {w}. \end{array}
\]

(We agree that, for all \(z \in \mathbf{C}\) , \(z^{1/2}\) is the square root of \(z\) whose amplitude belongs to \(\{0, \pi\}\) .) Denote by \(\mathrm{S}_{v,q}\) the \(\mathfrak{s}\) -module defined by \(\rho_{v,q}\) . We have \(\rho_{v,q}(\mathrm{C}) = 4q\) .

\begin{enumerate}
\def\labelenumi{\alph{enumi})}
\setcounter{enumi}{1}
\item
  If \(q \neq u^2 + u\) for all \(u \in v\) , \(S_{v,q}\) is simple.
\item
  Assume that \(2v \neq 0\) , and that q is of the form \(u^{2} + u\) , where \(u \in v, u \geq 0\) (which defines u uniquely). Let \(S_{v,q}^{-}\) (resp. \(S_{v,q}^{+}\) ) be the vector subspace of \(S_{v}\) generated by the \(e_{w}\) for \(w \leq u\) (resp. w \textgreater{} u). Then \(S_{v,q}^{-}\) and \(S_{v,q}^{+}\) are simple s-submodules of \(S_{v,q}\) .
\item
  Assume that 2v = 0, and that q is of the form \(u^{2} + u\) , where \(u \in v, u \geq 0\) (which defines u uniquely). Let \(S_{v,q}^{-}\) (resp. \(S_{v,q}^{0}, S_{v,q}^{+}\) ) be the vector subspace of \(S_{v}\) generated by the \(e_{w}\) for w \textless{} -u (resp. \(-u \leq w \leq u, w > u\) ). Then \(S_{v,q}^{-}, S_{v,q}^{0}\) and \(S_{v,q}^{+}\) are simple s-submodules of \(S_{v,q}\) .
\item
  Assume that \(v = -\frac{1}{2} + Z\) and \(q = -\frac{1}{4}\) . Let \(S_{-1/2,-1/4}^{-}\) (resp. \(S_{-1/2,-1/4}^{+}\) ) be the vector subspace of \(S_{-1/2}\) generated by the \(e_{w}\) for \(w \leq -\frac{1}{2}\) (resp. \(w > -\frac{1}{2}\) ). Then \(S_{-1/2,-1/4}^{-}\) and \(S_{-1/2,-1/4}^{+}\) are simple s-submodules of \(S_{-1/2,-1/4}\) .
\end{enumerate}

\(f)\) Denote by \(\rho_{v,q}^{\pm},\rho_{v,q}^{0}\) the representations corresponding to \(\mathrm{S}_{v,q}^{\pm},\mathrm{S}_{v,q}^{0}\) . In case \(b)\) , the elements \(x\in \mathfrak{s}\) such that \(\rho_{v,q}^{-}(x)\) admits an eigenvalue are those in \(\mathbf{CH}\) . In cases \(c)\) , \(d)\) , \(e)\) , the elements \(x\in \mathfrak{s}\) such that \(\rho_{v,q}^{-}(x)\) admits an eigenvalue are those in \(\mathbf{CH} + \mathbf{CX}_+\) ; if, in addition, \(x\) is nilpotent (and hence proportional to \(X_{+}\) ) and non-zero, \(\rho_{v,q}^{-}(x)\) admits 0 as its only eigenvalue, and this of multiplicity 1; on the other hand, if \(x\) is semi-simple, the underlying space of \(\rho_{v,q}^{-}\) has a basis consisting of eigenvectors of \(\rho_{v,q}^{-}(x)\) . There are analogous results for \(\rho_{v,q}^{+}\) , replacing \(X_{+}\) by \(X_{-}\) .

\(g)\) Let \(V\) be a simple \(\mathfrak{s}\) -module and \(\rho\) the corresponding representation. Then \(\rho(C)\) is a homothety (use Exerc. 11). Assume that \(\rho(C) = 4q\) . Show that, if \(\rho\) is \(H\) -diagonalizable, \(V\) is isomorphic to one of the modules \(S_{v,q}, S_{v,q}^{\pm}, S_{v,q}^{0}\) considered in \(b), c), d), e)\) . Moreover, \(\rho\) is \(H\) -diagonalizable if and only if \(\rho(H)\) admits an eigenvalue; it suffices that \(\rho(X_{+})\) admits the eigenvalue \(0\) .

¶14) The notations are those of the preceding exercise. Denote by \(B_{q}\) the quotient of \(U(\mathfrak{s})\) by the two-sided ideal generated by C - 4q, and denote by \(u \mapsto u^{\bullet}\) the canonical map \(U(\mathfrak{s}) \to B_{q}\) . Consider the representations of Exerc. 12 and 13 as representations of \(B_{q}\) .

\begin{enumerate}
\def\labelenumi{\alph{enumi})}
\tightlist
\item
  Every element of \(B_{q}\) can be expressed uniquely in the form
\end{enumerate}

\[
\sum_ {r \geq 0} X _ {+} ^ {\bullet r} p _ {r} (H ^ {\bullet}) + \sum_ {s > 0} q _ {s} (H ^ {\bullet}) X _ {-} ^ {\bullet s},
\]

where the \(p_{r}\) and the \(q_{s}\) are polynomials. If two elements of \(B_{q}\) generate the same left ideal, they are proportional.

\begin{enumerate}
\def\labelenumi{\alph{enumi})}
\setcounter{enumi}{1}
\item
  Consider case b) of Exerc. 13. Let a be a non-zero element of \(S_{v,q}\) . If a is an eigenvector of \(\rho_{v,q}(H)\) , with eigenvalue \(\lambda\) , the annihilator of a in \(B_{q}\) is the left ideal generated by \(H^{\bullet}-\lambda\) ; if a is not an eigenvector of \(\rho_{v,q}(H)\) , its annihilator is a non-monogenic left ideal.
\item
  Consider the case of the representations \(\rho_{v,q}^{\pm}\) of Exerc. 13. If a is a non-zero element of the space of such a representation, its annihilator in \(B_{q}\) is a non-monogenic left ideal.
\item
  Consider the representation \(\rho_{q}\) of Exerc. 12. The annihilator of \(e_{0}\) in \(B_{q}\) is generated by \(X_{+}^{\bullet}-1\) . If \(a\in W\) is not proportional to \(e_{0}\) , its annihilator is a non-monogenic left ideal.
\item
  Every automorphism of \(\mathfrak{s}\) extends to \(\mathrm{U}(\mathfrak{s})\) and defines by passage to the quotient an automorphism of \(\mathrm{B}_q\) ; let \(\mathrm{A}'\) be the subgroup of \(\mathrm{A} = \mathrm{Aut}(\mathrm{B}_q)\) thus obtained. Show that \(\mathrm{A}' \neq \mathrm{A}\) . (Let \(\varphi\) be the endomorphism \(x \mapsto [X_{+}^{\bullet 2}, x]\) of \(\mathrm{B}_q\) ; this automorphism is locally nilpotent, and \(e^{\varphi} \in \mathrm{A}, e^{\varphi} \notin \mathrm{A}'\) .)
\item
  The group A operates by transport of structure on the set of classes of simple representations of \(B_{q}\) . Let \(\pi_{1}\) (resp. \(\pi_{2},\pi_{3}\) ) be a representation of type \(\rho_{v,q}\) of Exerc. 13 b) (resp. of type \(\rho_{v,q}^{\pm}\) of Exerc. 13, resp. of type \(\rho_{q}\) of Exerc. 12). Then \(A\pi_{1},A\pi_{2}\) and \(A\pi_{3}\) are pairwise disjoint. (Use a), b), c), d). If \(\psi\in A\) is such that \(\psi\pi_{1}\) is of type \(\rho_{v,q}\) of Exerc. 13 b), then \(\psi\in A'\) (use a) and b)). Deduce that, if \(\psi\) is the automorphism \(e^{\varphi}\) of e), the representation \(\sigma=\pi_{1}\circ\psi\) has the following property: for all \(x\in s-\{0\}\) , \(\sigma(x)\) has no eigenvalue. \(^{8}\)
\end{enumerate}

\footnotetext[8]{For more details on Exerc. 12, 13 and 14, see: D. ARNAL and G. PINCZON, Sur les représentations algébriquement irréductibles de l'algèbre de Lie sl(2), J. Math. Phys., Vol. 15 (1974), pp. 350-359.}

\begin{enumerate}
\def\labelenumi{\arabic{enumi})}
\setcounter{enumi}{14}
\tightlist
\item
  Let \(\mathfrak{g}\) be a Lie algebra of dimension 3. Show the equivalence of the following conditions (cf.~Chap. I, §6, Exerc. 23).
\end{enumerate}

\begin{enumerate}
\def\labelenumi{(\roman{enumi})}
\item
  \(\mathfrak{g} = [\mathfrak{g},\mathfrak{g}]\) .
\item
  The Killing form of \(\mathfrak{g}\) is non-degenerate.
\item
  g is semi-simple.
\item
  g is simple.
\end{enumerate}

¶ 16) Let g be a simple Lie algebra of dimension 3, and let \(\Phi\) be its Killing form. Denote by \(\mathfrak{o}(\Phi)\) the orthogonal algebra of \(\Phi\) , i.e.~the subalgebra of \(\mathfrak{gl}(\mathfrak{g})\) consisting of the elements leaving \(\Phi\) invariant.

\begin{enumerate}
\def\labelenumi{\alph{enumi})}
\item
  Show that ad : \(\mathfrak{g} \to \mathfrak{o}(\varPhi)\) is an isomorphism.
\item
  Prove the equivalence of the following properties:
\end{enumerate}

\begin{enumerate}
\def\labelenumi{(\roman{enumi})}
\item
  \(\mathfrak{g}\) contains a non-zero isotropic vector (for \(\varPhi\) ).
\item
  \(\mathfrak{g}\) contains a non-zero nilpotent element.
\item
  g is isomorphic to s.
\end{enumerate}

\begin{enumerate}
\def\labelenumi{\alph{enumi})}
\setcounter{enumi}{2}
\item
  Show that there exists an extension \(k_{1}\) of \(k\) , of degree \(\leq 2\) , such that \(\mathfrak{g}_{(k_1)} = k_1 \otimes_k \mathfrak{g}\) is isomorphic to \(\mathfrak{s}_{(k_1)}\) .
\item
  Give the vector space \(\mathbf{A} = k\oplus \mathfrak{g}\) the unique algebra structure admitting 1 as unit element, and such that the product \(\mathfrak{g}\times \mathfrak{g}\to \mathrm{A}\) is given by the formula
\end{enumerate}

\[
x. y = \frac {1}{8} \Phi (x, y). 1 + \frac {1}{2} [ x, y ] \quad (x, y \in \mathfrak {g}).
\]

Show that A is a quaternion algebra over k, and that g is the Lie subalgebra of A consisting of the elements of trace zero (reduce, by extension of the base field, to the case where g = s, and show that A can then be identified with the matrix algebra \(\mathbf{M}_{2}(k)\) ).

Conversely, if D is a quaternion algebra over k, the elements of D of trace zero form a simple Lie algebra of dimension 3, and the corresponding algebra A can be identified with D.

\begin{enumerate}
\def\labelenumi{\alph{enumi})}
\setcounter{enumi}{4}
\tightlist
\item
  Prove the formulas
\end{enumerate}

\[
\begin{array}{c} \Phi (x, x) = - 8 \mathrm{Nrd} _ {\mathrm{A}} (x), \quad \Phi (x, y) = 4 \mathrm{Trd} _ {\mathrm{A}} (x y) \\ 2 [ x, [ y, z ] ] = \Phi (x, y) z - \Phi (x, z) y \qquad (x, y, z \in \mathfrak {g}). \end{array}
\]

\(f\) ) Show that the discriminant of \(\varPhi\) (with respect to any basis of \(\mathfrak{g}\) ) is of the form \(-2\lambda^2\) , with \(\lambda \in k^{*}\) .

\begin{enumerate}
\def\labelenumi{\alph{enumi})}
\setcounter{enumi}{6}
\tightlist
\item
  Let n be an integer \(\geq 0\) . Show that the g-module \(\mathbf{S}^{n}(\mathfrak{g})\) has a unique simple submodule of dimension \(2n + 1\) , and that this module is absolutely simple (reduce to the case where g = s, and use Exerc. 9).
\end{enumerate}

Show that g has an absolutely simple module of dimension 2n only if g is isomorphic to s, i.e.~only if A is isomorphic to \(\mathbf{M}_{2}(k)\) . (Let V be such a module. If \(n \geq 2\) , show by means of Exerc. 6 c) that the g-module \(V \otimes g\) has a unique absolutely simple submodule of dimension 2n - 2. Then reduce to the case n = 1, which is trivial.)

¶ 17) We retain the notations of the preceding exercise. Show that, for all \(n \geq 1\) , the algebra \(\mathrm{U}(\mathfrak{g})\) has a unique two-sided ideal \(m_{n}\) such that \(\mathrm{U}(\mathfrak{g}) / \mathfrak{m}_{n}\) is a central simple algebra of dimension \(n^{2}\) (extend scalars to reduce to the case in which g = s and show that \(m_{n}\) is then the kernel of the homomorphism \(\mathrm{U}(\mathfrak{g}) \to \mathrm{End}(\mathrm{V}(n-1))\) ). Every two-sided ideal of \(\mathrm{U}(\mathfrak{g})\) of finite codimension is of the form \(m_{n_{1}} \cap m_{n_{2}} \cap \cdots \cap m_{n_{h}}\) , where \(n_{1}, \ldots, n_{h}\) are distinct; its codimension is \(n_{1}^{2} + \cdots + n_{h}^{2}\) (apply the density theorem). The \(m_{n}\) are the only maximal two-sided ideals of \(\mathrm{U}(\mathfrak{g})\) of finite codimension.

Show that \(m_{2}\) is generated (as a two-sided ideal) by the elements \(x^{2}-\frac{1}{8}\Phi(x,x)\) ( \(x\in g\) ), and that the quotient \(\mathrm{U}(\mathfrak{g})/\mathfrak{m}_{2}\) can be identified with the quaternion algebra A of Exerc. 16.

When g is isomorphic to s, \(\mathrm{U}(\mathfrak{g})/\mathfrak{m}_{n}\) is isomorphic to \(\mathbf{M}_{n}(k)\) . Show that, when g is not isomorphic to s, \(\mathrm{U}(\mathfrak{g})/\mathfrak{m}_{n}\) is isomorphic to \(\mathbf{M}_{n}(k)\) if and only if n is odd (use Exerc. 16 g)). \(^{9}\)

\footnotetext[9]{It can be shown that, when n is even, \(\mathrm{U}(\mathfrak{g})/\mathfrak{m}_{n}\) is isomorphic to \(\mathbf{M}_{n/2}(\mathrm{A})\) .}

¶ 18) Assume that \(k = \mathbf{R}, \mathbf{C}\) or a field complete for a discrete valuation and with residue field of characteristic \(p \neq 0\) (for example a finite extension of the \(p\) -adic field \(\mathbf{Q}_p\) ).

\begin{enumerate}
\def\labelenumi{\alph{enumi})}
\tightlist
\item
  Let \(\mathfrak{n}\) be a nilpotent Lie algebra, N the Lie group obtained by giving \(\mathfrak{n}\) the Hausdorff law (Chap. III, §9, no. 5), and \(\rho : \mathfrak{n} \to \mathfrak{gl}(\mathrm{E})\) a linear representation of \(\mathfrak{n}\) on a finite dimensional vector space E. Assume that \(\rho(x)\) is nilpotent for all \(x \in \mathfrak{n}\) , and put \(\pi(x) = \exp(\rho(x))\) . Show that \(\pi\) is the only homomorphism of Lie groups \(\varphi : \mathrm{N} \to \mathbf{GL}(\mathrm{E})\) such that \(\mathrm{L}(\varphi) = \rho\) .
\end{enumerate}

(When k = R or C, use the fact that N is connected. When k is ultrametric, show that the eigenvalues of \(\varphi(x)\) , \(x \in N\) , are equal to 1; for this, prove first of all that, if \(k'\) is a finite extension of k, and if \((\lambda_{1}, \ldots, \lambda_{n}, \ldots)\) is a sequence of elements of \(k'\) such that \(\lambda_{n} = \lambda_{n+1}^{p}\) for all n and \(\lambda_{1} = 1\) , then \(\lambda_{n} = 1\) for all n.)

\begin{enumerate}
\def\labelenumi{\alph{enumi})}
\setcounter{enumi}{1}
\tightlist
\item
  Let \(\rho : \mathfrak{s} \to \mathfrak{gl}(\mathrm{E})\) be a finite dimensional linear representation of \(\mathfrak{s}\) , and let \(\pi\) be the homomorphism from \(\mathbf{SL}(2, k)\) to \(\mathbf{GL}(\mathrm{E})\) compatible with \(\rho\) (no. 4). Show that \(\pi\) is the unique homomorphism of Lie groups \(\varphi : \mathbf{SL}(2, k) \to \mathbf{GL}(\mathrm{E})\) such that \(\mathrm{L}(\varphi) = \rho\) . (Use \(a\) ) to prove that \(\pi\) and \(\rho\) coincide on the \(\exp(n)\) , with \(n\) nilpotent in \(\mathfrak{s}\) , and remark that \(\mathbf{SL}(2, k)\) is generated by the \(\exp(n)\) .)
\end{enumerate}

